\documentclass[10pt,a4paper]{amsart}
\usepackage[margin=19mm]{geometry}
\usepackage{amsmath,amssymb,mathtools}
\usepackage[scr=rsfs,cal=euler]{mathalfa}
\usepackage{xfrac}
\usepackage[unicode,hidelinks,bookmarks=false]{hyperref}
\newcommand{\ie}{i.e.\ }
\newcommand{\cf}{cf.\ }
\newcommand{\resp}{respectively\ }
\newcommand{\Subsection}{\subsection}
\providecommand{\nobreakdash}{\nobreak\hbox{-}}
\setlength{\emergencystretch}{2em}
\multlinegap0pt
\usepackage{mathtools}
\usepackage{amssymb}
\usepackage{bm}
\usepackage[shortlabels]{enumitem}
\usepackage{float}
\usepackage{xcolor}
\usepackage{tikz}
\usepackage[normalem]{ulem}
\newtheorem{theorem}{Theorem}
\newtheorem{lemma}{Lemma}
\newtheorem{claim}{Claim}
\newcommand{\cro}{{\mathsf{cr}}}
\newcommand{\dist}{\textnormal{dist}}
\newcommand{\mset}{{\mathcal M}}
\newcommand{\pset}{{\mathcal{P}}}
\newcommand{\set}[1]{\left\{ #1 \right\}}
\title{Paths and Intersections: Repelling Pairs}
\author{Yu Chen, Pavlo Pylyavskyy, Zihan Tan}
\date{Source: arXiv:2607.02883v2 (CC BY 4.0)}
\begin{document}
\maketitle

\noindent\emph{Source and licence.} Paths and Intersections: Repelling Pairs. Version arXiv:2607.02883v2 (CC BY 4.0), distributed by arXiv under the Creative Commons Attribution 4.0 International licence, http://creativecommons.org/licenses/by/4.0/, as recorded in the arXiv licence metadata for this exact version and shown on the abstract page https://arxiv.org/abs/2607.02883v2. Full licence notice: \texttt{licenses/arxiv-2607.02883-CC-BY-4.0.txt}.\par\smallskip

\noindent\textit{Editorial scope.} Counted unit: the article's lemma thm:minimum-realization (source0.tex lines 1501--1505). The article's complete original proof of it is delivered with it (source0.tex lines 1507--1585, one \texttt{proof} environment of 79 lines). No mathematics is written, repaired or re-derived: the statement and the proof are the article's own lines, in the article's own notation, and no retained line is substituted or re-flowed.  Retained uncounted as the source-local context the proof needs: lines 224--237; lines 1123--1126; lines 1488--1491.  Nothing was omitted from any retained span, so the package carries no omission record.  No retained line contains a citation, so the wrapper carries no bibliography and no bibliography entry.  No retained line contains a figure environment, an image-inclusion command or drawing code, so no image is delivered and the package declares no asset dependency.  Editorial formatting only, in addition: an AMS \texttt{amsart} preamble that restates the article's own declarations and macros used by the retained lines, namely the environments \texttt{theorem}, \texttt{lemma}, \texttt{claim} on separate counters, together with the standard packages those lines need.  The retained environments print in this document as Theorem 1, Lemma 1, Claim 1 in order of appearance, whereas the article prints the same items with the numbering of its own full document.  The complete native source of the article as distributed in the arXiv e-print for this exact version is bundled unchanged as \texttt{sources/204347/source0.tex}.
\begin{theorem}
\label{lem: 2-path condition}
For a Kalmanson metric $D$ on $T$, an Okamura-Seymour instance $(G,T)$ and a shortest path structure $\pset=\set{P_{t,t'}}_{t,t'\in T}$, the following are equivalent:
\begin{itemize}
    \item There exist nonnegative edge lengths $\set{\ell_e}_{e\in E(G)}$, such that
    for every pair $t,t'$ of terminals,
    \label{prop_2}
    \begin{itemize}
        \item $P_{t,t'}$ is indeed the shortest $t$-$t'$ path in $G$; and
        \item the length of $P_{t,t'}$ equals $D(t,t')$.
    \end{itemize}
    \item For every repelling pair $(t_1,t_2)$ and $(t'_1,t'_2)$, the paths $P_{t_1,t_2}$ and $P_{t'_1,t'_2}$ are disjoint.
\end{itemize}
\end{theorem}

\begin{theorem}
\label{thm:repelling-paths}
Given a metric $D$ on $T$, an OS instance $(G,T)$ admits a $D$-good shortest path structure iff for each repelling set $\mset$ and each chain $A$, the number of pairs in $\mset$ that cross $A$ is at most $|A|$. Moreover, such a $D$-good shortest path structure, if it exists, can be computed efficiently.
\end{theorem}

\begin{theorem}\label{thm:unique-template}
For every OS metric $D$, $\Phi(D)$ is the unique template with minimum crossing number among all templates satisfying $\dist_\Phi(x,y)\ge b_{x,y}$ for all $x,y\in X$.
That is, every other feasible template $\Phi\ne \Phi(D)$ satisfies that $\cro(\Phi)>\cro(\Phi(D))$.
\end{theorem}

\begin{lemma}\label{thm:minimum-realization}
For any OS metric $D$, the graph structures of all minimum realizations of $D$ is exactly the set of primal graphs of all arrangements of template $\Phi(D)$.
% with edge lengths assigned as in Theorem~\ref{lem: 2-path condition}.  
The number of edges in these minimum realizations is $\cro(\Phi(D))$.
\end{lemma}

\begin{proof}

Let $(G,T)$ be a minimum OS realization of $D$. Let $\Gamma$ be its
medial graph with medial template $\Phi$.  We first show a basic property of $\Gamma$.
\begin{claim}
If two medial chords
$\alpha$ and $\beta$ in $\Phi$ have non-alternating endpoints on the boundary, then they do not cross each other.    
\end{claim}
\begin{proof}
Suppose towards contradiction that they meet. Since their endpoints are non-alternating, they
meet an even positive number of times.  Hence two consecutive
intersection points of $\alpha$ and $\beta$ bound a lens.  Choose an
innermost such lens, its interior contains no medial vertex. After moving the crossing strands out of its interior (this operation does not change the number of edges in the primal graph) and obtaining an empty medial lens, in the
primal graph, such an empty medial lens can be simplified: either a degree-$2$ nonterminal, or a two-sided
face bounded by two parallel edges.  In the first case we replace the
two incident edges by one whose length is their sum; in the second
case we delete the longer of the two parallel edges.  Both operations
preserve all terminal distances and strictly decrease the number of
edges, contradicting the minimality of $G$.  Thus two medial chords
with non-alternating endpoints do not meet. 
\end{proof}  

The same lens argument
shows that two chords with alternating endpoints meet exactly once.
Consequently, $\Gamma$ is an arrangement of its template $\Phi$.
%
Now fix a boundary cut $x$-$y$, and let $L$ and $R$ be the two boundary
arcs determined by $x,y$.  A medial chord is called crossing if
it has one endpoint on $L$ and the other endpoint on $R$.  We prove the following claim. 
\begin{claim}
There is a simple arc $\gamma$ from $x$ to $y$ which meets precisely the
crossing chords, and meets each of them exactly once.    
\end{claim}
\begin{proof}
Indeed, if no such arc existed, then the chords whose two endpoints
lie on the same side of the $x$-$y$ cut would contain a topological
barrier between $x$ and $y$.  Equivalently, some connected component of
their union would meet both $L$ and $R$.  But every chord in this union
has both endpoints on $L$ or both endpoints on $R$.  Therefore any
component meeting both sides must contain an intersection between an
$L$-$L$ chord and an $R$-$R$ chord.  These two chords have
non-alternating endpoints, contradicting the reducedness proved above.
Hence such an arc $\gamma$ exists.  Since $\Gamma$ is an arrangement,
each crossing chord meets $\gamma$ exactly once, and no same-side
chord meets $\gamma$.
\end{proof}

Reading the primal and dual cells crossed by $\gamma$ gives an
$x$-$y$ chain $A$.  The primal cells in this sequence are precisely the
vertices of the chain.  Since $\gamma$ meets exactly the medial chords
crossing the $x$-$y$ cut, we obtain
$\operatorname{dist}_{\Phi}(x,y)=
    2|A|-\mathbf 1[x\in T]-\mathbf 1[y\in T]$.
The endpoint correction comes from the fact that a nonterminal boundary
point lies in a peripheral region, whereas a terminal endpoint lies at a
primal vertex.

Since $(G,T)$ realizes $D$, Theorem~\ref{lem: 2-path condition} implies that $G$ admits a
$D$-good shortest path structure.  Therefore Theorem~\ref{thm:repelling-paths} applies.  In
particular, every $x$-$y$ chain has length at least $a_{x,y}$, so the
chain $A$ above satisfies $|A|\ge a_{x,y}$.  Hence
$\operatorname{dist}_{\Phi}(x,y)
    \ge
    2a_{x,y}-\mathbf 1[x\in T]-\mathbf 1[y\in T]
    =
    b_{x,y}$.
Thus the medial template of every minimum realization is feasible.
By Theorem~\ref{thm:unique-template},
$\operatorname{cr}(\Phi)\ge \operatorname{cr}(\Phi(D)).
$
Thus every minimum realization has at least
$\operatorname{cr}(\Phi(D))$ edges.

On the other hand, take any arrangement of template $\Phi(D)$ and form the primal graph by the standard medial-to-primal construction. 
%bi-color the complementary regions so that the regions incident with terminal boundary arcs are primal, place a primal vertex in each primal region, and insert an edge for each crossing of two medial chords.  The boundary vertices corresponding to the terminal regions are the terminals $T$.  
Since $\dist_{\Phi(D)}(x,y)=b_{x,y}$, the minimum chain lengths in this primal graph are exactly the lower bounds $a_{x,y}$, with the endpoint correction.  Hence all chain inequalities of Theorem~\ref{thm:repelling-paths} hold, so the resulting graph admits a $D$-good shortest-path structure, and by Theorem~\ref{lem: 2-path condition} it can be assigned nonnegative edge lengths realizing $D$. The number of edges is $\cro(\Phi(D))$.

The uniqueness of the template follows from Theorem~\ref{thm:unique-template}.  The only freedom in the graph structure is the intersection pattern of chords with this fixed boundary pairing.
\end{proof}

\end{document}
